\xchapter{Conference IX. The First Conference of Abbot Isaac. On Prayer.}

\xsection{Chapter I. Introduction to the Conference.}

Introduction to the Conference.

\textsc{What} was promised in the second book of the Institutes\footnote{See the Institutes Book II. c. ix.} on continual and unceasing perseverance in prayer, shall be by the Lord’s help fulfilled by the Conferences of this Elder, whom we will now bring forward; viz., Abbot Isaac:\footnote{Isaac was, as we gathered from c. xxxi., a disciple of St. Antony, and is mentioned by Palladius Dial. de vita Chrysost. There are also a few stories of him in the Apophegmata Patrum (Migne, Vol. lxv. p. 223); and see the Dictionary of Christian Biography, Vol. iii. p. 294.} and when these have been propounded I think that I shall have satisfied the commands of Pope Castor of blessed memory, and your wishes, O blessed Pope Leontius and holy brother Helladius, and the length of the book in its earlier part may be excused, though, in spite of our endeavour not only to compress what had to be told into a brief discourse, but also to pass over very many points in silence, it has been extended to a greater length than we intended. For having commenced with a full discourse on various regulations which we have thought it well to curtail for the sake of brevity, at the close the blessed Isaac spoke these words.

\xsection{Chapter II. The words of Abbot Isaac on the nature of prayer.}

The words of Abbot Isaac on the nature of prayer.

\textsc{The} aim of every monk and the perfection of his heart tends to continual and unbroken perseverance in prayer, and, as far as it is allowed to human frailty, strives to acquire an immovable tranquillity of mind and a perpetual purity, for the sake of which we seek unweariedly and constantly to practise all bodily labours as well as contrition of spirit. And there is between these two a sort of reciprocal and inseparable union. For just as the crown of the building of all virtues is the perfection of prayer, so unless everything has been united and compacted by this as its crown, it cannot possibly continue strong and stable. For lasting and continual calmness in prayer, of which we are speaking, cannot be secured or consummated without them, so neither can those virtues which lay its foundations be fully gained without persistence in it. And so we shall not be able either to treat properly of the effect of prayer, or in a rapid discourse to penetrate to its main end, which is acquired by labouring at all virtues, unless first all those things which for its sake must be either rejected or secured, are singly enumerated and discussed, and, as the Parable in the gospel teaches,\footnote{Cf. S. Luke xiv. 28.} whatever concerns the building of that spiritual and most lofty tower, is reckoned up and carefully considered beforehand. But yet these things when prepared will be of no use nor allow the lofty height of perfection to be properly placed upon them unless a clearance of all faults be first undertaken, and the decayed and dead rubbish of the passions be dug up, and the strong foundations of simplicity and humility be laid on the solid and (so to speak) living soil of our breast, or rather on that rock of the gospel,\footnote{Cf. S. Luke vi. 48.} and by being built in this way this tower of spiritual virtues will rise, and be able to stand unmoved, and be raised to the utmost heights of heaven in full assurance of its stability. For if it rests on such foundations, then though heavy storms of passions break over it, though mighty torrents of persecutions beat against it like a battering ram, though a furious tempest of spiritual foes dash against it and attack it, yet not only will no ruin overtake it, but the onslaught will not injure it even in the slightest degree.

\xsection{Chapter III. How pure and sincere prayer can be gained.}

How pure and sincere prayer can be gained.

\textsc{And} therefore in order that prayer may be offered up with that earnestness and purity with which it ought to be, we must by all means observe these rules. First all anxiety about carnal things must be entirely got rid of; next we must leave no room for not merely the care but even the recollection of any business affairs, and in like manner also must lay aside all backbitings, vain and incessant chattering, and buffoonery; anger above all and disturbing moroseness must be entirely destroyed, and the deadly taint of carnal lust and covetousness be torn up by the roots. And so when these and such like faults which are also visible to the eyes of men, are entirely removed and cut off, and when such a purification and cleansing, as we spoke of, has first taken place, which is brought about by pure simplicity and innocence, then first there must be laid the secure foundations of a deep humility, which may be able to support a tower that shall reach the sky; and next the spiritual structure of the virtues must be built up upon them, and the soul kept free from all conversation and from roving thoughts that thus it may by little and little begin to rise to the contemplation of God and to spiritual insight. For whatever our mind has been thinking of before the hour of prayer, is sure to occur to us while we are praying through the activity of the memory. Wherefore what we want to find ourselves like while we are praying, that we ought to prepare ourselves to be before the time for prayer. For the mind in prayer is formed by its previous condition, and when we are applying ourselves to prayer the images of the same actions and words and thoughts will dance before our eyes, and make us either angry, as in our previous condition, or gloomy, or recall our former lust and business, or make us shake with foolish laughter (which I am ashamed to speak of) at some silly joke, or smile at some action, or fly back to our previous conversation. And therefore if we do not want anything to haunt us while we are praying, we should be careful before our prayer, to exclude it from the shrine of our heart, that we may thus fulfill the Apostle’s injunction: “Pray without ceasing;” and: “In every place lifting up holy hands without wrath or disputing.”\footnote{1 Thess. v. 17; 1 Tim. ii. 8.} For otherwise we shall not be able to carry out that charge unless our mind, purified from all stains of sin, and given over to virtue as to its natural good, feed on the continual contemplation of Almighty God.

\xsection{Chapter IV. Of the lightness of the soul which may be compared to a wing or feather.}

Of the lightness of the soul which may be compared to a wing or feather.

\textsc{For} the nature of the soul is not inaptly compared to a very fine feather or very light wing, which, if it has not been damaged or affected by being spoilt by any moisture falling on it from without, is borne aloft almost naturally to the heights of heaven by the lightness of its nature, and the aid of the slightest breath: but if it is weighted by any moisture falling upon it and penetrating into it, it will not only not be carried away by its natural lightness into any aerial flights but will actually be borne down to the depths of earth by the weight of the moisture it has received. So also our soul, if it is not weighted with faults that touch it, and the cares of this world, or damaged by the moisture of injurious lusts, will be raised as it were by the natural blessing of its own purity and borne aloft to the heights by the light breath of spiritual meditation; and leaving things low and earthly will be transported to those that are heavenly and invisible. Wherefore we are well warned by the Lord’s command: “Take heed that your hearts be not weighed down by surfeiting and drunkenness and the cares of this world.”\footnote{S. Luke xxi. 34.} And therefore if we want our prayers to reach not only the sky, but what is beyond the sky, let us be careful to reduce our soul, purged from all earthly faults and purified from every stain, to its natural lightness, that so our prayer may rise to God unchecked by the weight of any sin.

\xsection{Chapter V. Of the ways in which our soul is weighed down.}

Of the ways in which our soul is weighed down.

\textsc{But} we should notice the ways in which the Lord points out that the soul is weighed down: for He did not mention adultery, or fornication, or murder, or blasphemy, or rapine, which everybody knows to be deadly and damnable, but surfeiting and drunkenness, and the cares or anxieties of this world: which men of this world are so far from avoiding or considering damnable that actually some who (I am ashamed to say) call themselves monks entangle themselves in these very occupations as if they were harmless or useful. And though these three things, when literally given way to weigh down the soul, and separate it from God, and bear it down to things earthly, yet it is very easy to avoid them, especially for us who are separated by so great a distance from all converse with this world, and who do not on any occasion have anything to do with those visible cares and drunkenness and surfeiting. But there is another surfeiting which is no less dangerous, and a spiritual drunkenness which it is harder to avoid, and a care and anxiety of this world, which often ensnares us even after the perfect renunciation of all our goods, and abstinence from wine and all feastings and even when we are living in solitude—and of such the prophet says: “Awake, ye that are drunk but not with wine;”\footnote{Joel i. 5.} and another: “Be astonished and wonder and stagger: be drunk and not with wine: be moved, but not with drunkenness.”\footnote{Is. xxix. 9.} And of this drunkenness the wine must consequently be what the prophet calls “the fury of dragons”: and from what root the wine comes you may hear: “From the vineyard of Sodom,” he says, “is their vine, and their branches from Gomorrha.” Would you also know about the fruit of that vine and the seed of that branch? “Their grape is a grape of gall, theirs is a cluster of bitterness”\footnote{Deut. xxxii. 32, 33.} for unless we are altogether cleansed from all faults and abstaining from the surfeit of all passions, our heart will without drunkenness from wine and excess of any feasting be weighed down by a drunkenness and surfeiting that is still more dangerous. For that worldly cares can sometimes fall on us who mix with no actions of this world, is clearly shown according to the rule of the Elders, who have laid down that anything which goes beyond the necessities of daily food, and the unavoidable needs of the flesh, belongs to worldly cares and anxieties, as for example if, when a job bringing in a penny would satisfy the needs of our body, we try to extend it by a longer toil and work in order to get twopence or threepence; and when a covering of two tunics would be enough for our use both by night and day, we manage to become the owners of three or four, or when a hut containing one or two cells would be sufficient, in the pride of worldly ambition and greatness we build four or five cells, and these splendidly decorated, and larger than our needs required, thus showing the passion of worldly lusts whenever we can.

\xsection{Chapter VI. Of the vision which a certain Elder saw concerning the restless work of a brother.}

Of the vision which a certain Elder saw concerning the restless work of a brother.

\textsc{And} that this is not done without the prompting of devils we are taught by the surest proofs, for when one very highly esteemed Elder was passing by the cell of a certain brother who was suffering from this mental disease of which we have spoken, as he was restlessly toiling in his daily occupations in building and repairing what was unnecessary, he watched him from a distance breaking a very hard stone with a heavy hammer, and saw a certain Ethiopian standing over him and together with him striking the blows of the hammer with joined and clasped hands, and urging him on with fiery incitements to diligence in the work: and so he stood still for a long while in astonishment at the force of the fierce demon and the deceitfulness of such an illusion. For when the brother was worn out and tired and wanted to rest and put an end to his toil, he was stimulated by the spirit’s prompting and urged on to resume his hammer again and not to cease from devoting himself to the work which he had begun, so that being unweariedly supported by his incitements he did not feel the harm that so great labour was doing him. At last then the old man, disgusted at such a horrid mystification by a demon, turned aside to the brother’s cell and saluted him, and asked “what work is it, brother, that you are doing?” and he replied: “We are working at this awfully hard stone, and we can hardly break it at all.” Whereupon the Elder replied: “You were right in saying ‘\emph{we} can,’ for you were not alone, when you were striking it, but there was another with you whom you did not see, who was standing over you not so much to help you as urge you on with all his force.” And thus the fact that the disease of worldly vanity has not got hold of our hearts, will be proved by no mere abstinence from those affairs which even if we want to engage in, we cannot carry out, nor by the despising of those matters which if we pursued them would make us remarkable in the front rank among spiritual persons as well as among worldly men, but only when we reject with inflexible firmness of mind whatever ministers to our power and seems to be veiled in a show of right. And in reality these things which seem trivial and of no consequence, and which we see to be permitted indifferently by those who belong to our calling, none the less by their character affect the soul than those more important things, which according to their condition usually intoxicate the senses of worldly people and which do not allow\footnote{\emph{Sinentes}, though the reading of almost all \textsc{mss.} must be an error either of the author or of a copyist for \emph{sinentia}.} a monk to lay aside earthly impurities and aspire to God, on whom his attention should ever be fixed; for in his case even a slight separation from that highest good must be regarded as present death and most dangerous destruction. And when the soul has been established in such a peaceful condition, and has been freed from the meshes of all carnal desires, and the purpose of the heart has been steadily fixed on that which is the only highest good, he will then fulfil this Apostolic precept: “Pray without ceasing;” and: “in every place lifting up holy hands without wrath and disputing:”\footnote{1 Thess. v. 17; 1 Tim. ii. 8.} for when by this purity (if we can say so) the thoughts of the soul are engrossed, and are re-fashioned out of their earthly condition to bear a spiritual and angelic likeness, whatever it receives, whatever it takes in hand, whatever it does, the prayer will be perfectly pure and sincere.

\xsection{Chapter VII. A question how it is that it is harder work to preserve than to originate good thoughts.}

A question how it is that it is harder work to preserve than to originate good thoughts.

\textsc{Germanus}: If only we could keep as a lasting possession those spiritual thoughts in the same way and with the same ease with which we generally conceive their germs! for when they have been conceived in our hearts either through the recollection of the Scriptures or by the memory of some spiritual actions, or by gazing upon heavenly mysteries, they vanish all too soon and disappear by a sort of unnoticed flight. And when our soul has discovered some other occasions for spiritual emotions, different ones again crowd in upon us, and those which we had grasped are scattered, and lightly fly away so that the mind retaining no persistency, and keeping of its own power no firm hand over holy thoughts, must be thought, even when it does seem to retain them for a while, to have conceived them at random and not of set purpose. For how can we think that their rise should be ascribed to our own will, if they do not last and remain with us? But that we may not owing to the consideration of this question wander any further from the plan of the discourse we had commenced, or delay any longer the explanation promised of the nature of prayer, we will keep this for its own time, and ask to be informed at once of the character of prayer, especially as the blessed Apostle exhorts us at no time to cease from it, saying “Pray without ceasing.” And so we want to be taught first of its character, i.e., how prayer ought \emph{always} to be offered up, and then how we can secure this, whatever it is, and practise it without ceasing. For that it cannot be done by any light purpose of heart both daily experience and the explanation of four holiness show us, as you have laid it down that the aim of a monk, and the height of all perfection consist in the consummation of prayer.

\xsection{Chapter VIII. Of the different characters of prayer.}

Of the different characters of prayer.

\textsc{Isaac}: I imagine that all kinds of prayers cannot be grasped without great purity of heart and soul and the illumination of the Holy Spirit. For there are as many of them as there can be conditions and characters produced in one soul or rather in all souls. And so although we know that owing to our dulness of heart we cannot see all kinds of prayers, yet we will try to relate them in some order, as far as our slender experience enables us to succeed. For according to the degree of the purity to which each soul attains, and the character of the state in which it is sunk owing to what happens to it, or is by its own efforts renewing itself, its very prayers will each moment be altered: and therefore it is quite clear that no one can always offer up uniform prayers. For every one prays in one way when he is brisk, in another when he is oppressed with a weight of sadness or despair, in another when he is invigorated by spiritual achievements, in another when cast down by the burden of attacks, in another when he is asking pardon for his sins, in another when he asks to obtain grace or some virtue or else prays for the destruction of some sin, in another when he is pricked to the heart by the thought of hell and the fear of future judgment, in another when he is aglow with the hope and desire of good things to come, in another when he is taken up with affairs and dangers, in another when he is in peace and security, in another when he is enlightened by the revelation of heavenly mysteries, and in another when he is depressed by a sense of barrenness in virtues and dryness in feeling.

\xsection{Chapter IX. Of the fourfold nature of prayer.}

Of the fourfold nature of prayer.

\textsc{And} therefore, when we have laid this down with regard to the character of prayer, although not so fully as the importance of the subject requires, but as fully as the exigencies of time permit, and at any rate as our slender abilities admit, and our dulness of heart enables us,—a still greater difficulty now awaits us; viz., to expound one by one the different kinds of prayer, which the Apostle divides in a fourfold manner, when he says as follows: “I exhort therefore first of all that supplications, prayers, intercessions, thanksgivings be made.”\footnote{1 Tim. ii. 1.} And we cannot possibly doubt that this division was not idly made by the Apostle. And to begin with we must investigate what is meant by supplication, by prayer, by intercession, and by thanksgiving. Next we must inquire whether these four kinds are to be taken in hand by him who prays all at once, i.e., are they all to be joined together in every prayer,—or whether they are to be offered up in turns and one by one, as, for instance, ought at one time supplications, at another prayers, at another intercessions, and at another thanksgivings to be offered, or should one man present to God supplications, another prayers, another intercessions, another thanksgivings, in accordance with that measure of age, to which each soul is advancing by earnestness of purpose?

\xsection{Chapter X. Of the order of the different kinds laid down with regard to the character of prayer.}

Of the order of the different kinds laid down with regard to the character of prayer.

\textsc{And} so to begin with we must consider the actual force of the names and words, and discuss what is the difference between prayer and supplication and intercession; then in like manner we must investigate whether they are to be offered separately or all together; and in the third place must examine whether the particular order which is thus arranged by the Apostle’s authority has anything further to teach the hearer, or whether the distinction simply is to be taken, and it should be considered that they were arranged by him indifferently in such a way: a thing which seems to me utterly absurd. For one must not believe that the Holy Spirit uttered anything casually or without reason through the Apostle. And so we will, as the Lord grants us, consider them in the same order in which we began.

\xsection{Chapter XI. Of Supplications.}

Of Supplications.

“I \textsc{exhort} therefore first of all that supplications be made.” Supplication is an imploring or petition concerning sins, in which one who is sorry for his present or past deeds asks for pardon.

\xsection{Chapter XII. Of Prayer.}

Of Prayer.

\textsc{Prayers} are those by which we offer or vow something to God, what the Greeks call \textgreek{εὐκή}, i.e., a vow. For where we read in Greek \textgreek{ἰὰς ἐυκάς μου τῶ κυρίῶ ἀποδώσω}, in Latin we read: “I will pay my vows unto the Lord;”\footnote{Ps. cxv. 4 (cxvi. 14).} where according to the exact force of the words it may be thus represented: “I will pay my prayers unto the Lord.” And this which we find in Ecclesiastes: “If thou vowest a vow unto the Lord do not delay to pay it,” is written in Greek likewise: \textgreek{ἐάν ἐύξῃ ἐυχὴν τῶ κυρίῶ,} i.e., “If thou prayest a prayer unto the Lord, do not delay to pay it,”\footnote{Eccl. v. 3.} which will be fulfilled in this way by each one of us. We pray, when we renounce this world and promise that being dead to all worldly actions and the life of this world we will serve the Lord with full purpose of heart. We pray when we promise that despising secular honours and scorning earthly riches we will cleave to the Lord in all sorrow of heart and humility of spirit. We pray when we promise that we will ever maintain the most perfect purity of body and steadfast patience, or when we vow that we will utterly root out of our heart the roots of anger or of sorrow that worketh death. And if, enervated by sloth and returning to our former sins we fail to do this we shall be guilty as regards our prayers and vows, and these words will apply to us: “It is better not to vow, than to vow and not to pay,” which can be rendered in accordance with the Greek: “It is better for thee not to pray than to pray and not to pay.”\footnote{\emph{Ibid}. ver. 4.}

\xsection{Chapter XIII. Of Intercession.}

Of Intercession.

\textsc{In} the third place stand intercessions, which we are wont to offer up for others also, while we are filled with fervour of spirit, making request either for those dear to us or for the peace of the whole world, and to use the Apostle’s own phrase, we pray “for all men, for kings and all that are in authority.”\footnote{1 Tim. ii. 1, 2.}

\xsection{Chapter XIV. Of Thanksgiving.}

Of Thanksgiving.

\textsc{Then} in the fourth place there stand thanksgivings which the mind in ineffable transports offers up to God, either when it recalls God’s past benefits or when it contemplates His present ones, or when it looks forward to those great ones in the future which God has prepared for them that love Him. And with this purpose too sometimes we are wont to pour forth richer prayers, while, as we gaze with pure eyes on those rewards of the saints which are laid up in store hereafter, our spirit is stimulated to offer up unspeakable thanks to God with boundless joy.

\xsection{Chapter XV. Whether these four kinds of prayers are necessary for everyone to offer all at once or separately and in turns.}

Whether these four kinds of prayers are necessary for everyone to offer all at once or separately and in turns.

\textsc{And} of these four kinds, although sometimes occasions arise for richer and fuller prayers (for from the class of supplications which arises from sorrow for sin, and from the kind of prayer which flows from confidence in our offerings and the performance of our vows in accordance with a pure conscience, and from the intercession which proceeds from fervour of love, and from the thanksgiving which is born of the consideration of God’s blessings and His greatness and goodness, we know that oftentimes there proceed most fervent and ardent prayers so that it is clear that all these kinds of prayer of which we have spoken are found to be useful and needful for all men, so that in one and the same man his changing feelings will give utterance to pure and fervent petitions now of supplications, now of prayers, now of intercessions) yet the first seems to belong more especially to beginners, who are still troubled by the stings and recollection of their sins; the second to those who have already attained some loftiness of mind in their spiritual progress and the quest of virtue; the third to those who fulfil the completion of their vows by their works, and are so stimulated to intercede for others also through the consideration of their weakness, and the earnestness of their love; the fourth to those who have already torn from their hearts the guilty thorns of conscience, and thus being now free from care can contemplate with a pure mind the beneficence of God and His compassions, which He has either granted in the past, or is giving in the present, or preparing for the future, and thus are borne onward with fervent hearts to that ardent prayer which cannot be embraced or expressed by the mouth of men. Sometimes however the mind which is advancing to that perfect state of purity and which is already beginning to be established in it, will take in all these at one and the same time, and like some incomprehensible and all-devouring flame, dart through them all and offer up to God inexpressible prayers of the purest force, which the Spirit Itself, intervening with groanings that cannot be uttered, while we ourselves understand not, pours forth to God, grasping at that hour and ineffably pouring forth in its supplications things so great that they cannot be uttered with the mouth nor even at any other time be recollected by the mind. And thence it comes that in whatever degree any one stands, he is found sometimes to offer up pure and devout prayers; as even in that first and lowly station which has to do with the recollection of future judgment, he who still remains under the punishment of terror and the fear of judgment is so smitten with sorrow for the time being that he is filled with no less keenness of spirit from the richness of his supplications than he who through the purity of his heart gazes on and considers the blessings of God and is overcome with ineffable joy and delight. For, as the Lord Himself says, he begins to love the more, who knows that he has been forgiven the more.\footnote{Cf. S. Luke vii. 47.}

\xsection{Chapter XVI. Of the kinds of prayer to which we ought to direct ourselves.}

Of the kinds of prayer to which we ought to direct ourselves.

\textsc{Yet} we ought by advancing in life and attaining to virtue to aim rather at those kinds of prayer which are poured forth either from the contemplation of the good things to come or from fervour of love, or which at least, to speak more humbly and in accordance with the measure of beginners, arise for the acquirement of some virtue or the extinction of some fault. For otherwise we shall not possibly attain to those sublimer kinds of supplication of which we spoke, unless our mind has been little by little and by degrees raised through the regular course of those intercessions.

\xsection{Chapter XVII. How the four kinds of supplication were originated by the Lord.}

How the four kinds of supplication were originated by the Lord.

\textsc{These} four kinds of supplication the Lord Himself by His own example vouchsafed to originate for us, so that in this too He might fulfil that which was said of Him: “which Jesus began both to do and to teach.”\footnote{Acts i. 1.} For He made use of the class of \emph{supplication} when He said: “Father, if it be possible, let this cup pass from me;” or this which is chanted in His Person in the Psalm: “My God, My God, look upon Me, why hast Thou forsaken me,”\footnote{S. Matt. xxvi. 39; Ps. xxi. (xxii.) 2.} and others like it. It is \emph{prayer} where He says: “I have magnified Thee upon the earth, I have finished the work which Thou gavest Me to do,” and this: “And for their sakes I sanctify Myself that they also may be sanctified in the truth.”\footnote{S. John xvii. 4, 19.} It is \emph{intercession} when He says: “Father, those Whom Thou hast given me, I will that they also may be with Me that they may see My glory which Thou hast given Me;” or at any rate when He says: “Father, forgive them for they know not what they do.”\footnote{\emph{Ib}. 24; S. Luke xxiii. 34.} It is \emph{thanksgiving} when He says: “I confess to Thee, Father, Lord of heaven and earth, that Thou hast hid these things from the wise and prudent, and hast revealed them unto babes. Even so, Father, for so it seemed good in Thy sight:” or at least when He says: “Father, I thank Thee that Thou hast heard Me. But I knew that Thou hearest Me always.”\footnote{S. Matt. xi. 25, 26; S. John xi. 41, 42.} But though our Lord made a distinction between these four kinds of prayers as to be offered separately and one by one according to the scheme which we know of, yet that they can all be embraced in a perfect prayer at one and the same time He showed by His own example in that prayer which at the close of S. John’s gospel we read that He offered up with such fulness. From the words of which (as it is too long to repeat it all) the careful inquirer can discover by the order of the passage that this is so. And the Apostle also in his Epistle to the Philippians has expressed the same meaning, by putting these four kinds of prayers in a slightly different order, and has shown that they ought sometimes to be offered together in the fervour of a single prayer, saying as follows: “But in everything by prayer and supplication with thanksgiving let your requests be made known unto God.”\footnote{Phil. iv. 6.} And by this he wanted us especially to understand that in prayer and supplication thanksgiving ought to be mingled with our requests.

\xsection{Chapter XVIII. Of the Lord's Prayer.}

Of the Lord’s Prayer.

\textsc{And} so there follows after these different kinds of supplication a still more sublime and exalted condition which is brought about by the contemplation of God alone and by fervent love, by which the mind, transporting and flinging itself into love for Him, addresses God most familiarly as its own Father with a piety of its own. And that we ought earnestly to seek after this condition the formula of the Lord’s prayer teaches us, saying “Our Father.” When then we confess with our own mouths that the God and Lord of the universe is our Father, we profess forthwith that we have been called from our condition as slaves to the adoption of sons, adding next “Which art in heaven,” that, by shunning with the utmost horror all lingering in this present life, which we pass upon this earth as a pilgrimage, and what separates us by a great distance from our Father, we may the rather hasten with all eagerness to that country where we confess that our Father dwells, and may not allow anything of this kind, which would make us unworthy of this our profession and the dignity of an adoption of this kind, and so deprive us as a disgrace to our Father’s inheritance, and make us incur the wrath of His justice and severity. To which state and condition of sonship when we have advanced, we shall forthwith be inflamed with the piety which belongs to good sons, so that we shall bend all our energies to the advance not of our own profit, but of our Father’s glory, saying to Him: “Hallowed be Thy name,” testifying that our desire and our joy is His glory, becoming imitators of Him who said: “He who speaketh of himself, seeketh his own glory. But He who seeks the glory of Him who sent Him, the same is true and there is no unrighteousness in Him.”\footnote{S. John vii. 18.} Finally the chosen vessel being filled with this feeling wished that he could be anathema from Christ\footnote{Cf. Rom. ix. 3.} if only the people belonging to Him might be increased and multiplied, and the salvation of the whole nation of Israel accrue to the glory of His Father; for with all assurance could he wish to die for Christ as he knew that no one perished for life. And again he says: “We rejoice when we are weak but ye are strong.”\footnote{2 Cor. xiii. 9.} And what wonder if the chosen vessel wished to be anathema from Christ for the sake of Christ’s glory and the conversion of His own brethren and the privilege of the nation, when the prophet Micah wished that he might be a liar and a stranger to the inspiration of the Holy Ghost, if only the people of the Jews might escape those plagues and the going forth into captivity which he had announced in his prophecy, saying: “Would that I were not a man that hath the Spirit, and that I rather spoke a lie;”\footnote{Micah ii. 11.}—to pass over that wish of the Lawgiver, who did not refuse to die together with his brethren who were doomed to death, saying: “I beseech Thee, O Lord; this people hath sinned a heinous sin; either forgive them this trespass, or if Thou do not, blot me out of Thy book which Thou hast written.”\footnote{Exod. xxxii. 31, 32.} But where it is said “Hallowed be Thy name,” it may also be very fairly taken in this way: “The hallowing of God is our perfection.” And so when we say to Him “Hallowed be Thy name” we say in other words, make us, O Father, such that we maybe able both to understand and take in what the hallowing of Thee is, or at any rate that Thou mayest be seen to be hallowed in our spiritual converse. And this is effectually fulfilled in our case when “men see our good works, and glorify our Father Which is in heaven.”\footnote{S. Matt. v. 16.}

\xsection{Chapter XIX. Of the clause “Thy kingdom come.”}

Of the clause “Thy kingdom come.”

\textsc{The} second petition of the pure heart desires that the kingdom of its Father may come at once; viz., either that whereby Christ reigns day by day in the saints (which comes to pass when the devil’s rule is cast out of our hearts by the destruction of foul sins, and God begins to hold sway over us by the sweet odour of virtues, and, fornication being overcome, charity reigns in our hearts together with tranquillity, when rage is conquered; and humility, when pride is trampled under foot) or else that which is promised in due time to all who are perfect, and to all the sons of God, when it will be said to them by Christ: “Come ye blessed of My Father, inherit the kingdom prepared for you from the foundation of the world;”\footnote{S. Matt. xxv. 34.} (as the heart) with fixed and steadfast gaze, so to speak, yearns and longs for it and says to Him “Thy kingdom come.” For it knows by the witness of its own conscience that when He shall appear, it will presently share His lot. For no guilty person would dare either to say or to wish for this, for no one would want to face the tribunal of the Judge, who knew that at His coming he would forthwith receive not the prize or reward of his merits but only punishment.

\xsection{Chapter XX. Of the clause “Thy will be done.”}

Of the clause “Thy will be done.”

\textsc{The} third petition is that of sons: “Thy will be done as in heaven so on earth.” There can now be no grander prayer than to wish that earthly things may be made equal with things heavenly: for what else is it to say “Thy will be done as in heaven so on earth,” than to ask that men may be like angels and that as God’s will is ever fulfilled by them in heaven, so also all those who are on earth may do not their own but His will? This too no one could say from the heart but only one who believed that God disposes for our good all things which are seen, whether fortunate or unfortunate, and that He is more careful and provident for our good and salvation than we ourselves are for ourselves. Or at any rate it may be taken in this way: The will of God is the salvation of all men, according to these words of the blessed Paul: “Who willeth all men to be saved and to come to the knowledge of the truth.”\footnote{1 Tim. ii. 4.} Of which will also the prophet Isaiah says in the Person of God the Father: “And all Thy will shall be done.”\footnote{Is. xlvi. 10.} When we say then “Thy will be done as in heaven so on earth,” we pray in other words for this; viz., that as those who are in heaven, so also may all those who dwell on earth be saved, O Father, by the knowledge of Thee.

\xsection{Chapter XXI. Of our supersubstantial or daily bread.}

Of our supersubstantial or daily bread.

\textsc{Next}: “Give us this day our bread which is \textgreek{ἐπιούσιον},” i.e., “supersubstantial,” which another Evangelist calls “daily.”\footnote{Here Cassian is relying entirely on Jerome’s revised text of the Latin, which has \emph{supersubstantialis} in S. Matt. vi. 11, as the rendering of \textgreek{ἐπιούσιος} but translates the same word by \emph{quotidianum} in the parallel passage in S. Luke xi. 3. It is curious that Cassian should have been thus misled, with his knowledge of Greek, as well as his acquaintance with the old Latin version which has \emph{quotidianum} in both gospels. Cf. Bishop Lightfoot “On a Fresh Revision of the New Testament,” p. 219.} The former indicates the quality of its nobility and substance, in virtue of which it is above all substances and the loftiness of its grandeur and holiness exceeds all creatures, while the latter intimates the purpose of its use and value. For where it says “daily” it shows that without it we cannot live a spiritual life for a single day. Where it says “today” it shows that it must be received daily and that yesterday’s supply of it is not enough, but at it must be given to us today also in like manner. And our daily need of it suggests to us that we ought at all times to offer up this prayer, because there is no day on which we have no need to strengthen the heart of our inner man, by eating and receiving it, although the expression used, “today” may be taken to apply to his present life, i.e., while we are living in this world supply us with this bread. For we know that it will be given to those who deserve it by Thee hereafter, but we ask that Thou wouldest grant it to us today, because unless it has been vouchsafed to a man to receive it in this life he will never be partaker of it in that.

\xsection{Chapter XXII. Of the clause: “Forgive us our debts, etc.”}

Of the clause: “Forgive us our debts, etc.”

“\textsc{And} forgive us our debts as we also forgive our debtors.” O unspeakable mercy of God, which has not only given us a form of prayer and taught us a system of life acceptable to Him, and by the requirements of the form given, in which He charged us always to pray, has torn up the roots of both anger and sorrow, but also gives to those who pray an opportunity and reveals to them a way by which they may move a merciful and kindly judgment of God to be pronounced over them and which somehow gives us a power by which we can moderate the sentence of our Judge, drawing Him to forgive our offences by the example of our forgiveness: when we say to Him: “Forgive us as we also forgive.” And so without anxiety and in confidence from this prayer a man may ask for pardon of his own offences, if he has been forgiving towards his own debtors, and not towards those of his Lord. For some of us, which is very bad, are inclined to show ourselves calm and most merciful in regard to those things which are done to God’s detriment, however great the crimes may be, but to be found most hard and inexorable exactors of debts to ourselves even in the case of the most trifling wrongs. Whoever then does not from his heart forgive his brother who has offended him, by this prayer calls down upon himself not forgiveness but condemnation, and by his own profession asks that he himself may be judged more severely, saying: Forgive me as I also have forgiven. And if he is repaid according to his own request, what else will follow but that he will be punished after his own example with implacable wrath and a sentence that cannot be remitted? And so if we want to be judged mercifully, we ought also to be merciful towards those who have sinned against us. For only so much will be remitted to us, as we have remitted to those who have injured us however spitefully. And some dreading this, when this prayer is chanted by all the people in church, silently omit this clause, for fear lest they may seem by their own utterance to bind themselves rather than to excuse themselves, as they do not understand that it is in vain that they try to offer these quibbles to the Judge of all men, who has willed to show us beforehand how He will judge His suppliants. For as He does not wish to be found harsh and inexorable towards them, He has marked out the manner of His judgment, that just as we desire to be judged by Him, so we should also judge our brethren, if they have wronged us in anything, for “he shall have judgment without mercy who hath shown no mercy.”\footnote{S. James ii. 13.}

\xsection{Chapter XXIII. Of the clause: “Lead us not into temptation.“}

Of the clause: “Lead us not into temptation.”

\textsc{Next} there follows: “And lead us not into temptation,” on which there arises no unimportant question, for if we pray that we may not be suffered to be tempted, how then will our power of endurance be proved, according to this text: “Every one who is not tempted is not proved;”\footnote{Ecclus. xxxiv. 11.} and again: “Blessed is the man that endureth temptation?”\footnote{S. James i. 12.} The clause then, “Lead us not into temptation,” does not mean this; viz., do not permit us ever to be tempted, but do not permit us when we fall into temptation to be overcome. For Job was tempted, but was not led into temptation. For he did not ascribe folly to God nor blasphemy, nor with impious mouth did he yield to that wish of the tempter toward which he was drawn. Abraham was tempted, Joseph was tempted, but neither of them was led into temptation for neither of them yielded his consent to the tempter. Next there follows: “But deliver us from evil,” i.e., do not suffer us to be tempted by the devil above that we are able, but “make with the temptation a way also of escape that we may be able to bear it.”\footnote{1 Cor. x. 13.}

\xsection{Chapter XXIV. How we ought not to ask for other things, except only those which are contained in the limits of the Lord's Prayer.}

How we ought not to ask for other things, except only those which are contained in the limits of the Lord’s Prayer.

\textsc{You} see then what is the method and form of prayer proposed to us by the Judge Himself, who is to be prayed to by it, a form in which there is contained no petition for riches, no thought of honours, no request for power and might, no mention of bodily health and of temporal life. For He who is the Author of Eternity would have men ask of Him nothing uncertain, nothing paltry, and nothing temporal. And so a man will offer the greatest insult to His Majesty and Bounty, if he leaves on one side these eternal petitions and chooses rather to ask of Him something transitory and uncertain; and will also incur the indignation rather than the propitiation of the Judge by the pettiness of his prayer.

\xsection{Chapter XXV. Of the character of the sublimer prayer.}

Of the character of the sublimer prayer.

\textsc{This} prayer then though it seems to contain all the fulness of perfection, as being what was originated and appointed by the Lord’s own authority, yet lifts those to whom it belongs to that still higher condition of which we spoke above, and carries them on by a loftier stage to that ardent prayer which is known and tried by but very few, and which to speak more truly is ineffable; which transcends all human thoughts, and is distinguished, I will not say by any sound of the voice, but by no movement of the tongue, or utterance of words, but which the mind enlightened by the infusion of that heavenly light describes in no human and confined language, but pours forth richly as from copious fountain in an accumulation of thoughts, and ineffably utters to God, expressing in the shortest possible space of time such great things that the mind when it returns to its usual condition cannot easily utter or relate. And this condition our Lord also similarly prefigured by the form of those supplications which, when he retired alone in the mountain He is said to have poured forth in silence, and when being in an agony of prayer He shed forth even drops of blood, as an example of a purpose which it is hard to imitate.

\xsection{Chapter XXVI. Of the different causes of conviction.}

Of the different causes of conviction.

\textsc{But} who is able, with whatever experience he may be endowed, to give a sufficient account of the varieties and reasons and grounds of conviction, by which the mind is inflamed and set on fire and incited to pure and most fervent prayers? And of these we will now by way of specimen set forth a few, as far as we can by God’s enlightenment recollect them. For sometimes a verse of any one of the Psalms gives us an occasion of ardent prayer while we are singing. Sometimes the harmonious modulation of a brother’s voice stirs up the minds of dullards to intense supplication. We know also that the enunciation and the reverence of the chanter adds greatly to the fervour of those who stand by. Moreover the exhortation of a perfect man, and a spiritual conference has often raised the affections of those present to the richest prayer. We know too that by the death of a brother or some one dear to us, we are no less carried away to full conviction. The recollection also of our coldness and carelessness has sometimes aroused in us a healthful fervour of spirit. And in this way no one can doubt that numberless opportunities are not wanting, by which through God’s grace the coldness and sleepiness of our minds can be shaken off.

\xsection{Chapter XXVII. Of the different sorts of conviction.}

Of the different sorts of conviction.

\textsc{But} how and in what way those very convictions are produced from the inmost recesses of the soul it is no less difficult to trace out. For often through some inexpressible delight and keenness of spirit the fruit of a most salutary conviction arises so that it actually breaks forth into shouts owing to the greatness of its incontrollable joy; and the delight of the heart and greatness of exultation makes itself heard even in the cell of a neighbour. But sometimes the mind hides itself in complete silence within the secrets of a profound quiet, so that the amazement of a sudden illumination chokes all sounds of words and the overawed spirit either keeps all its feelings to itself or loses\footnote{Petschenig’s text reads “amittat.” v. l. emittat.} them and pours forth its desires to God with groanings that cannot be uttered. But sometimes it is filled with such overwhelming conviction and grief that it cannot express it except by floods of tears.

\xsection{Chapter XXVIII. A question about the fact that a plentiful supply of tears is not in our own power.}

A question about the fact that a plentiful supply of tears is not in our own power.

\textsc{Germanus}: My own poor self indeed is not altogether ignorant of this feeling of conviction. For often when tears arise at the recollection of my faults, I have been by the Lord’s visitation so refreshed by this ineffable joy which you describe that the greatness of the joy has assured me that I ought not to despair of their forgiveness. Than which state of mind I think there is nothing more sublime if only it could be recalled at our own will. For sometimes when I am desirous to stir myself up with all my power to the same conviction and tears, and place before my eyes all my faults and sins, I am unable to bring back that copiousness of tears, and so my eyes are dry and hard like some hardest flint, so that not a single tear trickles from them. And so in proportion as I congratulate myself on that copiousness of tears, just so do I mourn that I cannot bring it back again whenever I wish.

\xsection{Chapter XXIX. The answer on the varieties of conviction which spring from tears.}

The answer on the varieties of conviction which spring from tears.

\textsc{Isaac}: Not every kind of shedding of tears is produced by one feeling or one virtue. For in one way does that weeping originate which is caused by the pricks of our sins smiting our heart, of which we read: “I have laboured in my groanings, every night I will wash my bed; I will water my couch with my tears.”\footnote{Ps. vi. 7.} And again: “Let tears run down like a torrent day and night: give thyself no rest, and let not the apple of thine eye cease.”\footnote{Lam. ii. 18.} In another, that which arises from the contemplation of eternal good things and the desire of that future glory, owing to which even richer well-springs of tears burst forth from uncontrollable delights and boundless exultation, while our soul is athirst for the mighty Living God, saying, “When shall I come and appear before the presence of God? My tears have been my meat day and night,”\footnote{Ps. xii. (xliii.) 3, 4.} declaring with daily crying and lamentation: “Woe is me that my sojourning is prolonged;” and: “Too long hath my soul been a sojourner.”\footnote{Ps. cix. (cxix.) 5, 6.} In another way do the tears flow forth, which without any conscience of deadly sin, yet still proceed from the fear of hell and the recollection of that terrible judgment, with the terror of which the prophet was smitten and prayed to God, saying: “Enter not into judgment with Thy servant, for in Thy sight shall no man living be justified.”\footnote{Ps. cxlii. (cxliii.) 2.} There is too another kind of tears, which are caused not by knowledge of one’s self but by the hardness and sins of others; whereby Samuel is described as having wept for Saul, and both the Lord in the gospel and Jeremiah in former days for the city of Jerusalem, the latter thus saying: “Oh, that my head were water and mine eyes a fountain of tears! And I will weep day and night for the slain of the daughter of my people.”\footnote{Jer. ix. 1.} Or also such as were those tears of which we hear in the hundred and first Psalm: “For I have eaten ashes for my bread, and mingled my cup with weeping.”\footnote{Ps. ci. (cii.) 10.} And these were certainty not caused by the same feeling as those which arise in the sixth Psalm from the person of the penitent, but were due to the anxieties of this life and its distresses and losses, by which the righteous who are living in this world are oppressed. And this is clearly shown not only by the words of the Psalm itself, but also by its title, which runs as follows in the character of that poor person of whom it is said in the gospel that “blessed are the poor in spirit, for theirs is the kingdom of heaven:”\footnote{S. Matt. v. 3.} “A prayer of the poor when he was in distress and poured forth his prayer to God.”\footnote{Ps. ci. (cii.) 1.}

\xsection{Chapter XXX. How tears ought not to be squeezed out, when they do not flow spontaneously.}

How tears ought not to be squeezed out, when they do not flow spontaneously.

\textsc{From} these tears those are vastly different which are squeezed out from dry eyes while the heart is hard: and although we cannot believe that these are altogether fruitless (for the attempt to shed them is made with a good intention, especially by those who have not yet been able to attain to perfect knowledge or to be thoroughly cleansed from the stains of past or present sins), yet certainly the flow of tears ought not to be thus forced out by those who have already advanced to the love of virtue, nor should the weeping of the outward man be with great labour attempted, as even if it is produced it will never attain the rich copiousness of spontaneous tears. For it will rather cast down the soul of the suppliant by his endeavours, and humiliate him, and plunge him in human affairs and draw him away from the celestial heights, wherein the awed mind of one who prays should be steadfastly fixed, and will force it to relax its hold on its prayers and grow sick from barren and forced tears.

\xsection{Chapter XXXI. The opinion of Abbot Antony on the condition of prayer.}

The opinion of Abbot Antony on the condition of prayer.

\textsc{And} that you may see the character of true prayer I will give you not my own opinion but that of the blessed Antony: whom we have known sometimes to have been so persistent in prayer that often as he was praying in a transport of mind, when the sunrise began to appear, we have heard him in the fervour of his spirit declaiming: Why do you hinder me, O sun, who art arising for this very purpose; viz., to withdraw me from the brightness of this true light? And his also is this heavenly and more than human utterance on the end of prayer: That is not, said he, a perfect prayer, wherein a monk understands himself and the words which he prays. And if we too, as far as our slender ability allows, may venture to add anything to this splendid utterance, we will bring forward the marks of prayer which are heard from the Lord, as far as we have tried them.

\xsection{Chapter XXXII. Of the proof of prayer being heard.}

Of the proof of prayer being heard.

\textsc{When}, while we are praying, no hesitation intervenes and breaks down the confidence of our petition by a sort of despair, but we feel that by pouring forth our prayer we have obtained what we are asking for, we have no doubt that our prayers have effectually reached God. For so far will one be heard and obtain an answer, as he believes that he is regarded by God, and that God can grant it. For this saying of our Lord cannot be retracted: “Whatsoever ye ask when ye pray, believe that you shall receive, and they shall come to you.”\footnote{S. Mark xi. 24.}

\xsection{Chapter XXXIII. An objection that the confidence of being thus heard as described belongs only to saints.}

An objection that the confidence of being thus heard as described belongs only to saints.

\textsc{Germanus}: We certainly believe that this confidence of being heard flows from purity of conscience, but for us, whose heart is still smitten by the pricks of sins, how can we have it, as we have no merits to plead for us, whereby we might confidently presume that our prayers would be heard?

\xsection{Chapter XXXIV. Answer on the different reasons for prayer being heard.}

Answer on the different reasons for prayer being heard.

\textsc{Isaac}: That there are different reasons for prayer being heard in accordance with the varied and changing condition of souls the words of the gospels and of the prophets teach us. For you have the fruits of an answer pointed out by our Lord’s words in the case of the agreement of two persons; as it is said: “If two of you shall agree upon earth touching anything for which they shall ask, it shall be done for them of my Father which is in heaven.”\footnote{S. Matt. xviii. 19.} You have another in the fulness of faith, which is compared to a grain of mustard-seed. “For,” He says, “if you have faith as a grain of mustard seed, ye shall say unto this mountain: Be thou removed, and it shall be removed; and nothing shall be impossible to you.”\footnote{S. Matt. xvii. 19.} You have it in continuance in prayer, which the Lord’s words call, by reason of unwearied perseverance in petitioning, importunity: “For, verily, I say unto you that if not because of his friendship, yet because of his importunity he will rise and give him as much as he needs.”\footnote{S. Luke xi. 8.} You have it in the fruits of almsgiving: “Shut up alms in the heart of the poor and it shall pray for thee in the time of tribulation.”\footnote{Ecclus. xxix. 15.} You have it in the purifying of life and in works of mercy, as it is said: “Loose the bands of wickedness, undo the bundles that oppress;” and after a few words in which the barrenness of an unfruitful fast is rebuked, “then,” he says, “thou shalt call and the Lord shall hear thee; thou shalt cry, and He shall say, Here am I.”\footnote{Is. lviii. 6, 9.} Sometimes also excess of trouble causes it to be heard, as it is said: “When I was in trouble I called unto the Lord, and He heard me;”\footnote{Ps. cxix. (cxx.) 1.} and again: “Afflict not the stranger for if he crieth unto Me, I will hear him, for I am merciful.”\footnote{Exod. xxii. 21, 27.} You see then in how many ways the gift of an answer may be obtained, so that no one need be crushed by the despair of his conscience for securing those things which are salutary and eternal. For if in contemplating our wretchedness I admit that we are utterly destitute of all those virtues which we mentioned above, and that we have neither that laudable agreement of two persons, nor that faith which is compared to a grain of mustard seed, nor those works of piety which the prophet describes, surely we cannot be without that importunity which He supplies to all who desire it, owing to which alone the Lord promises that He will give whatever He has been prayed to give. And therefore we ought without unbelieving hesitation to persevere, and not to have the least doubt that by continuing in them we shall obtain all those things which we have asked according to the mind of God. For the Lord, in His desire to grant what is heavenly and eternal, urges us to constrain Him as it were by our importunity, as He not only does not despise or reject the importunate, but actually welcomes and praises them, and most graciously promises to grant whatever they have perseveringly hoped for; saying, “Ask and ye shall receive: seek and ye shall find: knock and it shall be opened unto you. For every one that asketh receiveth, and he that seeketh findeth, and to him that knocketh it shall be opened;”\footnote{S. Luke xi. 9, 10.} and again: “All things whatsoever ye shall ask in prayer believing ye shall receive, and nothing shall be impossible to you.”\footnote{S. Matt. xxi. 22; xvii. 20.} And therefore even if all the grounds for being heard which we have mentioned are altogether wanting, at any rate the earnestness of importunity may animate us, as this is placed in the power of any one who wills without the difficulties of any merits or labours. But let not any suppliant doubt that he certainly will not be heard, so long as he doubts whether he is heard. But that this also shall be sought from the Lord unweariedly, we are taught by the example of the blessed Daniel, as, though he was heard from the first day on which he began to pray, he only obtained the result of his petition after one and twenty days.\footnote{Cf. Dan. x. 2 \emph{sq}.} Wherefore we also ought not to grow slack in the earnestness of the prayers we have begun, if we fancy that the answer comes but slowly, for fear lest perhaps the gift of the answer be in God’s providence delayed, or the angel, who was to bring the Divine blessing to us, may when he comes forth from the Presence of the Almighty be hindered by the resistance of the devil, as it is certain that he cannot transmit and bring to us the desired boon, if he finds that we slack off from the earnestness of the petition made. And this would certainly have happened to the above mentioned prophet unless he had with incomparable steadfastness prolonged and persevered in his prayers until the twenty-first day. Let us then not be at all cast down by despair from the confidence of this faith of ours, even when we fancy that we are far from having obtained what we prayed for, and let us not have any doubts about the Lord’s promise where He says: “All things, whatsoever ye shall ask in prayer believing, ye shall receive.”\footnote{S. Matt. xxi. 22.} For it is well for us to consider this saying of the blessed Evangelist John, by which the ambiguity of this question is clearly solved: “This is,” he says, “the confidence which we have in Him, that whatsoever we ask according to His will, He heareth us.”\footnote{1 John v. 16.} He bids us then have a full and undoubting confidence of the answer only in those things which are not for our own advantage or for temporal comforts, but are in conformity to the Lord’s will. And we are also taught to put this into our prayers by the Lord’s Prayer, where we say “Thy will be done,”—\emph{Thine} not ours. For if we also remember these words of the Apostle that “we know not what to pray for as we ought”\footnote{Rom. viii. 26.} we shall see that we sometimes ask for things opposed to our salvation and that we are most providentially refused our requests by Him who sees what is good for us with greater right and truth than we can. And it is clear that this also happened to the teacher of the Gentiles when he prayed that the messenger of Satan who had been for his good allowed by the Lord’s will to buffet him, might be removed, saying: “For which I besought the Lord thrice that he might depart from me. And He said unto me, My grace is sufficient for thee, for strength is made perfect in weakness.”\footnote{2 Cor. xii. 8, 9.} And this feeling even our Lord expressed when He prayed in the character\footnote{\emph{Ex persona hominis assumpti}. The language is scarcely accurate, but it must be remembered that the Conferences were written before the rise of the Nestorian heresy had shown the need for exactness of expression on the subject of the Incarnation. Compare the note on “Against Nestorius,” Book III. c. iii.} of man which He had taken, that He might give us a form of prayer as other things also by His example; saying thus: “Father, if it be possible, let this cup pass from me: nevertheless not as I will but as Thou wilt,”\footnote{S. Matt. xxvi. 39.} though certainly His will was not discordant with His Father’s will, “For He had come to save what was lost and to give His life a ransom for many;”\footnote{S. Matt. xviii. 11; xx. 28.} as He Himself says: “No man taketh my life from Me, but I lay it down of Myself. I have power to lay it down and I have power to take it again.”\footnote{S. John x. 18.} In which character there is in the thirty-ninth Psalm the following sung by the blessed David, of the Unity of will which He ever maintained with the Father: “To do Thy will: O My God, I am willing.”\footnote{Ps. xxxix. (xl.) 9.} For even if we read of the Father: “For God so loved the world that He gave His only begotten Son,”\footnote{1 John iii. 16.} we find none the less of the Son: “Who gave Himself for our sins.”\footnote{Gal. i. 4.} And as it is said of the One: “Who spared not His own Son, but gave Him for all of us,”\footnote{Rom. viii. 32.} so it is written of the other: “He was offered because He Himself willed it.”\footnote{Is. liii. 7. (Lat.)} And it is shown that the will of the Father and of the Son is in all things one, so that even in the actual mystery of the Lord’s resurrection we are taught that there was no discord of operation. For just as the blessed Apostle declares that the Father brought about the resurrection of His body, saying: “And God the Father, who raised Him from the dead,”\footnote{Gal. i. 1.} so also the Son testifies that He Himself will raise again the Temple of His body, saying: “Destroy this temple, and in three days I will raise it up again.”\footnote{S. John ii. 19.} And therefore we being instructed by all these examples of our Lord which have been enumerated ought to end our supplications also with the same prayer, and always to subjoin this clause to all our petitions: “Nevertheless not as I will, but as Thou wilt.”\footnote{S. Matt. xxvi. 39.} But it is clear enough that one who does not\footnote{“Non” though wanting in most \textsc{mss.} must be read in the text.} pray with attention of mind cannot observe that threefold reverence\footnote{Reading “curvationis” with Petschenig: the text of Gazæus has “orationis.”} which is usually practised in the assemblies of the brethren at the close of service.

\xsection{Chapter XXXV. Of prayer to be offered within the chamber and with the door shut.}

Of prayer to be offered within the chamber and with the door shut.

\textsc{Before} all things however we ought most carefully to observe the Evangelic precept, which tells us to enter into our chamber and shut the door and pray to our Father, which may be fulfilled by us as follows: We pray within our chamber, when removing our hearts inwardly from the din of all thoughts and anxieties, we disclose our prayers in secret and in closest intercourse to the Lord. We pray with closed doors when with closed lips and complete silence we pray to the searcher not of words but of hearts. We pray in secret when from the heart and fervent mind we disclose our petitions to God alone, so that no hostile powers are even able to discover the character of our petition. Wherefore we should pray in complete silence, not only to avoid distracting the brethren standing near by our whispers or louder utterances, and disturbing the thoughts of those who are praying, but also that the purport of our petition may be concealed from our enemies who are especially on the watch against us while we are praying. For so we shall fulfil this injunction: “Keep the doors of thy mouth from her who sleepeth in thy bosom.”\footnote{Micah vii. 5.}

\xsection{Chapter XXXVI. Of the value of short and silent prayer.}

Of the value of short and silent prayer.

\textsc{Wherefore} we ought to pray often but briefly, lest if we are long about it our crafty foe may succeed in implanting something in our heart. For that is the true sacrifice, as “the sacrifice of God is a broken spirit.” This is the salutary offering, these are pure drink offerings, that is the “sacrifice of righteousness,” the “sacrifice of praise,” these are true and fat victims, “holocausts full of marrow,” which are offered by contrite and humble hearts, and which those who practise this control and fervour of spirit, of which we have spoken, with effectual power can sing: “Let my prayer be set forth in Thy sight as the incense: let the lifting up of my hands be an evening sacrifice.”\footnote{Ps. l. (li.) 19, 21; xlix. (l.) 23; lxv. (lxvi.) 15; cxl. (cxli.) 2.} But the approach of the right hour and of night warns us that we ought with fitting devotion to do this very thing, of which, as our slender ability allowed, we seem to have propounded a great deal, and to have prolonged our conference considerably, though we believe that we have discoursed very little when the magnificence and difficulty of the subject are taken into account.

With these words of the holy Isaac we were dazzled rather than satisfied, and after evening service had been held, rested our limbs for a short time, and intending at the first dawn again to return under promise of a fuller discussion departed, rejoicing over the acquisition of these precepts as well as over the assurance of his promises. Since we felt that though the excellence of prayer had been shown to us, still we had not yet understood from his discourse its nature, and the power by which continuance in it might be gained and kept.

\xchapter{Conference X. The Second Conference of Abbot Isaac. On Prayer.}

\xsection{Chapter I. Introduction.}

Introduction.

\textsc{Among} the sublime customs of the anchorites which by God’s help have been set forth although in plain and unadorned style, the course of our narration compels us to insert and find a place for something, which may seem so to speak to cause a blemish on a fair body: although I have no doubt that by it no small instruction on the image of Almighty God of which we read in Genesis will be conferred on some of the simpler sort, especially when the grounds are considered of a doctrine so important that men cannot be ignorant of it without terrible blasphemy and serious harm to the Catholic faith.

\xsection{Chapter II. Of the custom which is kept up in the Province of Egypt for signifying the time of Easter.}

Of the custom which is kept up in the Province of Egypt for signifying the time of Easter.

\textsc{In} the country of Egypt this custom is by ancient tradition observed that—when Epiphany is past, which the priests of that province regard as the time, both of our Lord’s baptism and also of His birth in the flesh, and so celebrate the commemoration of either mystery not separately as in the Western provinces but on the single festival of this day,\footnote{The observance of Epiphany can be traced back in the Christian Church to the second century, and, as Cassian tells us here, in the East (in which its observance apparently originated) it was in the first instance a double festival commemorating both the Nativity and the Baptism of our Lord. From the East its observance passed over to the West, where however the Nativity was already observed as a separate festival, and hence the \emph{special} reference of Epiphany was somewhat altered, and the manifestation to the Magi was coupled with that at the Baptism: hence the plural \emph{Epiphaniorum dies}. Meanwhile, as the West adopted the observance of this festival from the East, so the East followed the West in observing a separate feast of the Nativity. Cassian’s words show us that when he wrote the two festivals were both observed separately in the West, though apparently \emph{not} yet (to the best of his belief) in the East, but the language of a homily by S. Chrysostom (Vol. ii. p. 354 Ed. Montfaucon) delivered in \textsc{a.d.} 386 shows that the separation of the two festivals had already begun at Antioch, and all the evidence goes to show that “the Western plan was being gradually adopted in the period which we may roughly define as the last quarter of the 4th and the first quarter of the 5th century.” Dictionary of Christian Antiquities, Vol. i. p. 361. See further Origines du Culte Chrétien, par L’Abbé Duchesne, p. 247 \emph{sq}.}—letters are sent from the Bishop of Alexandria through all the Churches of Egypt, by which the beginning of Lent, and the day of Easter are pointed out not only in all the cities but also in all the monasteries.\footnote{The “Festal letters” (\textgreek{ἑοταστικαὶ ἐπιστολαί}, Euseb. VII. xx., xxi.) were delivered by the Bishop of Alexandria as Homilies, and then put into the form of an Epistle and sent round to all the churches of Egypt; and, according to some late writers, to the Bishops of all the principal sees, in accordance with a decision of the Council of Nicæa, in order to inform them of the right day on which Easter should be celebrated. Cassian here speaks of them as sent immediately after Epiphany, and this was certainly the time at which the announcement of the date of Easter was made in the West shortly after his day (so the Council of Orleans, Canon i., \textsc{a.d.} 541); that of Braga \textsc{a.d.} 572, Canon ix., and that of Auxerre \textsc{a.d.} 572, Canon ii.), but there is ample evidence in the Festal letters both of S. Athanasius and of S. Cyril that at Alexandria the homilies were preached on the previous Easter, and it is difficult to resist the inference that Cassian’s memory is here at fault as to the exact time at which the incident related really occurred, and that he is transferring to Egypt the custom with which he was familiar in the West, assigning to the festival of Epiphany what really must have taken place at Easter.} In accordance then with this custom, a very few days after the previous conference had been held with Abbot Isaac, there arrived the festal letters of Theophilus\footnote{Theophilus succeeded Timothy as Bishop of Alexandria in the summer of 385. The festal letters of which Cassian here speaks were issued by him in the year 399.} the Bishop of the aforesaid city, in which together with the announcement of Easter he considered as well the foolish heresy of the Anthropomorphites\footnote{The Anthropomorphite heresy, into which the monks of Egypt had fallen, “supposed that God possesses eyes, a face, and hands and other members of a bodily organization.” It arose from taking too literally those passages of the Old Testament in which God is spoken of in human terms, out of condescension to man’s limited powers of grasping the Divine nature and appears historically to have been a recoil from the allegorism of Origen and others of the Alexandrian school. The Festal letter of Theophilus in which he condemned these views, and maintained the incorporeal nature of God is no longer extant, but is alluded to also by Sozomen, H. E. VIII. xi., where an account is given of the Origenistic controversy of which it was the occasion, and out of which Theophilus came so badly. On the heresy see also Epiphanius, Hær. lxx.; Augustine. Hær. l. and lxxvi.; and Theodoret, H. E. IV. x.} at great length, and abundantly refuted it. And this was received by almost all the body of monks residing in the whole province of Egypt with such bitterness owing to their simplicity and error, that the greater part of the Elders decreed that on the contrary the aforesaid Bishop ought to be abhorred by the whole body of the brethren as tainted with heresy of the worst kind, because he seemed to impugn the teaching of holy Scripture by the denial that Almighty God was formed in the fashion of a human figure, though Scripture teaches with perfect clearness that Adam was created in His image. Lastly this letter was rejected also by those who were living in the desert of Scete and who excelled all who were in the monasteries of Egypt, in perfection and in knowledge, so that except Abbot Paphnutius the presbyter of our congregation, not one of the other presbyters, who presided over the other three churches in the same desert, would suffer it to be even read or repeated at all in their meetings.

\xsection{Chapter III. Of Abbot Sarapion and the heresy of the Anthropomorphites into which he fell in the error of simplicity.}

Of Abbot Sarapion and the heresy of the Anthropomorphites into which he fell in the error of simplicity.

\textsc{Among} those then who were caught by this mistaken notion was one named Sarapion, a man of long-standing strictness of life, and one who was altogether perfect in actual discipline, whose ignorance with regard to the view of the doctrine first mentioned was so far a stumbling block to all who held the true faith, as he himself outstripped almost all the monks both in the merits of his life and in the length of time (he had been there). And when this man could not be brought back to the way of the right faith by many exhortations of the holy presbyter Paphnutius, because this view seemed to him a novelty, and one that was not ever known to or handed down by his predecessors, it chanced that a certain deacon, a man of very great learning, named Photinus, arrived from the region of Cappadocia with the desire of visiting the brethren living in the same desert: whom the blessed Paphnutius received with the warmest welcome, and in order to confirm the faith which had been stated in the letters of the aforesaid Bishop, placed him in the midst and asked him before all the brethren how the Catholic Churches throughout the East interpreted the passage in Genesis where it says “Let us make man after our image and likeness.”\footnote{Gen. i. 26.} And when he explained that the image and likeness of God was taken by all the leaders of the churches not according to the base sound of the letters, but spiritually, and supported this very fully and by many passages of Scripture, and showed that nothing of this sort could happen to that infinite and incomprehensible and invisible glory, so that it could be comprised in a human form and likeness, since its nature is incorporeal and uncompounded and simple, and what can neither be apprehended by the eyes nor conceived by the mind, at length the old man was shaken by the numerous and very weighty assertions of this most learned man, and was drawn to the faith of the Catholic tradition. And when both Abbot Paphnutius and all of us were filled with intense delight at his adhesion, for this reason; viz., that the Lord had not permitted a man of such age and crowned with such virtues, and one who erred only from ignorance and rustic simplicity, to wander from the path of the right faith up to the very last, and when we arose to give thanks, and were all together offering up our prayers to the Lord, the old man was so bewildered in mind during his prayer because he felt that the Anthropomorphic image of the Godhead which he used to set before himself in prayer, was banished from his heart, that on a sudden he burst into a flood of bitter tears and continual sobs, and cast himself down on the ground and exclaimed with strong groanings: “Alas! wretched man that I am! they have taken away my God from me, and I have now none to lay hold of; and whom to worship and address I know not.” By which scene we were terribly disturbed, and moreover with the effect of the former Conference still remaining in our hearts, we returned to Abbot Isaac, whom when we saw close at hand, we addressed with these words.

\xsection{Chapter IV. Of our return to Abbot Isaac and question concerning the error into which the aforesaid old man had fallen.}

Of our return to Abbot Isaac and question concerning the error into which the aforesaid old man had fallen.

\textsc{Although} even besides the fresh matter which has lately arisen, our delight in the former conference which was held on the character of prayer would summon us to postpone everything else and return to your holiness, yet this grievous error of Abbot Sarapion, conceived, as we fancy, by the craft of most vile demons, adds somewhat to this desire of ours. For it is no small despair by which we are cast down when we consider that through the fault of this ignorance he has not only utterly lost all those labours which he has performed in so praiseworthy a manner for fifty years in this desert, but has also incurred the risk of eternal death. And so we want first to know why and wherefore so grievous an error has crept into him. And next we should like to be taught how we can arrive at that condition in prayer, of which you discoursed some time back not only fully but splendidly. For that admirable Conference has had this effect upon us, that it has only dazzled our minds and has not shown us how to perform or secure it.

\xsection{Chapter V. The answer on the heresy described above.}

The answer on the heresy described above.

\textsc{Isaac}: We need not be surprised that a really simple man who had never received any instruction on the substance and nature of the Godhead could still be entangled and deceived by an error of simplicity and the habit of a longstanding mistake, and (to speak more truly) continue in the original error which is brought about, not as you suppose by a new illusion of the demons, but by the ignorance of the ancient heathen world, while in accordance with the custom of that erroneous notion, by which they used to worship devils formed in the figure of men, they even now think that the incomprehensible and ineffable glory of the true Deity should be worshipped under the limitations of some figure, as they believe that they can grasp and hold nothing if they have not some image set before them, which they can continually address while they are at their devotions, and which they can carry about in their mind and have always fixed before their eyes. And against this mistake of theirs this text may be used: “And they changed the glory of the incorruptible God into the likeness of the image of corruptible man.”\footnote{Rom. i. 23.} Jeremiah also says: “My people have changed their glory for an idol.\footnote{Jer. ii. 11.} Which error although by this its origin, of which we have spoken, it is engrained in the notions of some, yet none the less is it contracted in the hearts also of those who have never been stained with the superstition of the heathen world, under the colour of this passage where it is said “Let us make man after our image and our likeness,”\footnote{Gen. i. 26.} ignorance and simplicity being its authors, so that actually there has arisen owing to this hateful interpretation a heresy called that of the Anthropomorphites, which maintains with obstinate perverseness that the infinite and simple substance of the Godhead is fashioned in our lineaments and human configuration. Which however any one who has been taught the Catholic doctrine will abhor as heathenish blasphemy, and so will arrive at that perfectly pure condition in prayer which will not only not connect with its prayers any figure of the Godhead or bodily lineaments (which it is a sin even to speak of), but will not even allow in itself even the memory of a name, or the appearance of an action, or an outline of any character.

\xsection{Chapter VI. Of the reasons why Jesus Christ appears to each one of us either in His humility or in His glorified condition.}

Of the reasons why Jesus Christ appears to each one of us either in His humility or in His glorified condition.

\textsc{For} according to the measure of its purity, as I said in the former Conference, each mind is both raised and moulded in its prayers if it forsakes the consideration of earthly and material things so far as the condition of its purity may carry it forward, and enable it with the inner eyes of the soul to see Jesus either still in His humility and in the flesh, or glorified and coming in the glory of His Majesty: for those cannot see Jesus coming in His Kingdom who are still kept back in a sort of state of Jewish weakness, and cannot say with the Apostle: “And if we have known Christ after the flesh, yet now we know Him so no more;”\footnote{2 Cor. v. 16.} but only those can look with purest eyes on His Godhead, who rise with Him from low and earthly works and thoughts and go apart in the lofty mountain of solitude which is free from the disturbance of all earthly thoughts and troubles, and secure from the interference of all sins, and being exalted by pure faith and the heights of virtue reveals the glory of His Face and the image of His splendour to those who are able to look on Him with pure eyes of the soul. But Jesus is seen as well by those who live in towns and villages and hamlets, i.e., who are occupied in practical affairs and works, but not with the same brightness with which He appeared to those who can go up with Him into the aforesaid mount of virtues, i.e., Peter, James, and John. For so in solitude He appeared to Moses and spoke with Elias. And as our Lord wished to establish this and to leave us examples of perfect purity, although He Himself, the very fount of inviolable sanctity, had no need of external help and the assistance of solitude in order to secure it (for the fulness of purity could not be soiled by any stain from crowds, nor could He be contaminated by intercourse with men, who cleanses and sanctifies all things that are polluted) yet still He retired into the mountain alone to pray, thus teaching us by the example of His retirement that if we too wish to approach God with a pure and spotless affection of heart, we should also retire from all the disturbance and confusion of crowds, so that while still living in the body we may manage in some degree to adapt ourselves to some likeness of that bliss which is promised hereafter to the saints, and that “God may be” to us “all in all.”\footnote{1 Cor. xv. 28.}

\xsection{Chapter VII. What constitutes our end and perfect bliss.}

What constitutes our end and perfect bliss.

\textsc{For} then will be perfectly fulfilled in our case that prayer of our Saviour in which He prayed for His disciples to the Father saying “that the love wherewith Thou lovedst Me may be in them and they in us;” and again: “that they all may be one as Thou, Father, in Me and I in Thee, that they also may be one in us,”\footnote{S. John xvii. 26, 21.} when that perfect love of God, wherewith “He first loved us”\footnote{1 John iv. 16.} has passed into the feelings of our heart as well, by the fulfilment of this prayer of the Lord which we believe cannot possibly be ineffectual. And this will come to pass when God shall be all our love, and every desire and wish and effort, every thought of ours, and all our life and words and breath, and that unity which already exists between the Father and the Son, and the Son and the Father, has been shed abroad in our hearts and minds, so that as He loves us with a pure and unfeigned and indissoluble love, so we also may be joined to Him by a lasting and inseparable affection, since we are so united to Him that whatever we breathe or think, or speak is God, since, as I say, we attain to that end of which we spoke before, which the same Lord in His prayer hopes may be fulfilled in us: “that they all may be one as we are one, I in them and Thou in Me, that they also may be made perfect in one;” and again: “Father, those whom Thou hast given Me, I will that where I am, they may also be with Me.”\footnote{S. John xvii. 22–24.} This then ought to be the destination of the solitary, this should be all his aim that it may be vouchsafed to him to possess even in the body an image of future bliss, and that he may begin in this world to have a foretaste of a sort of earnest of that celestial life and glory. This, I say, is the end of all perfection, that the mind purged from all carnal desires may daily be lifted towards spiritual things, until the whole life and all the thoughts of the heart become one continuous prayer.

\xsection{Chapter VIII. A question on the training in perfection by which we can arrive at perpetual recollection of God.}

A question on the training in perfection by which we can arrive at perpetual recollection of God.

\textsc{Germanus}: The extent of our bewilderment at our wondering awe at the former Conference, because of which we came back again, increases still more. For in proportion as by the incitements of this teaching we are fired with the desire of perfect bliss, so do we fall back into greater despair, as we know not how to seek or obtain training for such lofty heights. Wherefore we entreat that you will patiently allow us (for it must perhaps be set forth and unfolded with a good deal of talk) to explain what while sitting in the cell we had begun to revolve in a lengthy meditation, although we know that your holiness is not at all troubled by the infirmities of the weak, which even for this reason should be openly set forth, that what is out of place in them may receive correction. Our notion then is that the perfection of any art or system of training must begin with some simple rudiments, and grow accustomed first to somewhat easy and tender beginnings, so that being nourished and trained little by little by a sort of reasonable milk, it may grow up and so by degrees and step by step mount up from the lowest depths to the heights: and when by these means it has entered on the plainer principles and so to speak passed the gates of the entrance of the profession, it will consequently arrive without difficulty at the inmost shrine and lofty heights of perfection. For how could any boy manage to pronounce the simplest union of syllables unless he had first carefully learnt the letters of the alphabet? Or how can any one learn to read quickly, who is still unfit to connect together short and simple sentences? But by what means will one who is ill instructed in the science of grammar attain eloquence in rhetoric or the knowledge of philosophy? Wherefore for this highest learning also, by which we are taught even to cleave to God, I have no doubt that there are some foundations of the system, which must first be firmly laid and afterwards the towering heights of perfection may be placed and raised upon them. And we have a slight idea that these are its first principles; viz., that we should first learn by what meditations God may be grasped and contemplated, and next that we should manage to keep a very firm hold of this topic whatever it is which we do not doubt is the height of all perfection. And therefore we want you to show us some material for this recollection, by which we may conceive and ever keep the idea of God in the mind, so that by always keeping it before our eyes, when we find that we have dropped away from Him, we may at once be able to recover ourselves and return thither and may succeed in laying hold of it again without any delay from wandering around the subject and searching for it. For it happens that when we have wandered away from our spiritual speculations and have come back to ourselves as if waking from a deadly sleep, and, being thoroughly roused, look for the subject matter, by which we may be able to revive that spiritual recollection which has been destroyed, we are hindered by the delay of the actual search before we find it, and are once more drawn aside from our endeavour, and before the spiritual insight is brought about, the purpose of heart which had been conceived, has disappeared. And this trouble is certain to happen to us for this reason because we do not keep something special firmly set before our eyes like some principle to which the wandering thoughts may be recalled after many digressions and varied excursions; and, if I may use the expression, after long storms enter a quiet haven. And so it comes to pass that as the mind is constantly hindered by this want of knowledge and difficulty, and is always tossed about vaguely, and as if intoxicated, among various matters, and cannot even retain firm hold for any length of time of anything spiritual which has occurred to it by chance rather than of set purpose: while, as it is always receiving one thing after another, it does not notice either their beginning and origin or even their end.

\xsection{Chapter IX. The answer on the efficacy of understanding, which is gained by experience.}

The answer on the efficacy of understanding, which is gained by experience.

\textsc{Isaac}: Your minute and subtle inquiry affords an indication of purity being very nearly reached. For no one would be able even to make inquiries on these matters, I will not say to look within and discriminate,—except one who had been urged to sound the depths of such questions by careful and effectual diligence of mind, and watchful anxiety, and one whom the constant aim after a well controlled life had taught by practical experience to attempt the entrance to this purity and to knock at its doors. And therefore as I see you, I will not say, standing before the doors of that true prayer of which we have been speaking, but touching its inner chambers and inward parts as it were with the hands of experience, and already laying hold of some parts of it, I do not think that I shall find any difficulty in introducing you now within what I may call its hall, for you to roam about its recesses, as the Lord may direct; nor do I think that you will be hindered from investigating what is to be shown you by any obstacles or difficulties. For he is next door to understanding who carefully recognizes what he ought to ask about, nor is he far from knowledge, who begins to understand how ignorant he is. And therefore I am not afraid of the charge of betraying secrets, and of levity, if I divulge what when speaking in my former discourse on the perfection of prayer I had kept back from discussing, as I think that its force was to be explained to us who are occupied with this subject and interest even without the aid of my words, by the grace of God.

\xsection{Chapter X. Of the method of continual prayer.}

Of the method of continual prayer.

\textsc{Wherefore} in accordance with that system, which you admirably compared to teaching children (who can only take in the first lessons on the alphabet and recognize the shapes of the letters, and trace out their characters with a steady hand if they have, by means of some copies and shapes carefully impressed on wax, got accustomed to express their figures, by constantly looking at them and imitating them daily), we must give you also the form of this spiritual contemplation, on which you may always fix your gaze with the utmost steadiness, and both learn to consider it to your profit in unbroken continuance, and also manage by the practice of it and by meditation to climb to a still loftier insight. This formula then shall be proposed to you of this system, which you want, and of prayer, which every monk in his progress towards continual recollection of God, is accustomed to ponder, ceaselessly revolving it in his heart, having got rid of all kinds of other thoughts; for he cannot possibly keep his hold over it unless he has freed himself from all bodily cares and anxieties. And as this was delivered to us by a few of those who were left of the oldest fathers, so it is only divulged by us to a very few and to those who are really keen. And so for keeping up continual recollection of God this pious formula is to be ever set before you. “O God, make speed to save me: O Lord, make haste to help me,”\footnote{Ps. lxix. (lxx.) 2. It is not impropable that this chapter suggested to S. Benedict the use of these words as the opening versicle of the hour services, a position which it has ever since occupied in the West. See the rule of S. Benedict, cc. ix., xvii., and xviii.} for this verse has not unreasonably been picked out from the whole of Scripture for this purpose. For it embraces all the feelings which can be implanted in human nature, and can be fitly and satisfactorily adapted to every condition, and all assaults. Since it contains an invocation of God against every danger, it contains humble and pious confession, it contains the watchfulness of anxiety and continual fear, it contains the thought of one’s own weakness, confidence in the answer, and the assurance of a present and ever ready help. For one who is constantly calling on his protector, is certain that He is always at hand. It contains the glow of love and charity, it contains a view of the plots, and a dread of the enemies, from which one, who sees himself day and night hemmed in by them, confesses that he cannot be set free without the aid of his defender. This verse is an impregnable wall for all who are labouring under the attacks of demons, as well as impenetrable coat of mail and a strong shield. It does not suffer those who are in a state of moroseness and anxiety of mind, or depressed by sadness or all kinds of thoughts to despair of saving remedies, as it shows that He, who is invoked, is ever looking on at our struggles and is not far from His suppliants. It warns us whose lot is spiritual success and delight of heart that we ought not to be at all elated or puffed up by our happy condition, which it assures us cannot last without God as our protector, while it implores Him not only always but even speedily to help us. This verse, I say, will be found helpful and useful to every one of us in whatever condition we may be. For one who always and in all matters wants to be helped, shows that he needs the assistance of God not only in sorrowful or hard matters but also equally in prosperous and happy ones, that he may be delivered from the one and also made to continue in the other, as he knows that in both of them human weakness is unable to endure without His assistance. I am affected by the passion of gluttony. I ask for food of which the desert knows nothing, and in the squalid desert there are wafted to me odours of royal dainties and I find that even against my will I am drawn to long for them. I must at once say: “O God, make speed to save me: O Lord, make haste to help me.” I am incited to anticipate the hour fixed for supper, or I am trying with great sorrow of heart to keep to the limits of the right and regular meagre fare. I must cry out with groans: “O God, make speed to save me: O Lord, make haste to help me.” Weakness of the stomach hinders me when wanting severer fasts, on account of the assaults of the flesh, or dryness of the belly and constipation frightens me. In order that effect may be given to my wishes, or else that the fire of carnal lust may be quenched without the remedy of a stricter fast, I must pray: “O God, make speed to save me: O Lord, make haste to help me.” When I come to supper, at the bidding of the proper hour I loathe taking food and am prevented from eating anything to satisfy the requirements of nature: I must cry with a sigh: “O God, make speed to save me: O Lord, make haste to help me.” When I want for the sake of steadfastness of heart to apply myself to reading a headache interferes and stops me, and at the third hour sleep glues my head to the sacred page, and I am forced either to overstep or to anticipate the time assigned to rest; and finally an overpowering desire to sleep forces me to cut short the canonical rule for service in the Psalms: in the same way I must cry out: “O God, make speed to save me: O Lord, make haste to help me.” Sleep is withdrawn from my eyes, and for many nights I find myself wearied out with sleeplessness caused by the devil, and all repose and rest by night is kept away from my eyelids; I must sigh and pray: “O God, make speed to save me: O Lord, make haste to help me.” While I am still in the midst of a struggle with sin suddenly an irritation of the flesh affects me and tries by a pleasant sensation to draw me to consent while in my sleep. In order that a raging fire from without may not burn up the fragrant blossoms of chastity, I must cry out: “O God, make speed to save me: O Lord, make haste to help me.” I feel that the incentive to lust is removed, and that the heat of passion has died away in my members: In order that this good condition acquired, or rather that this grace of God may continue still longer or forever with me, I must earnestly say: “O God, make speed to save me: O Lord, make haste to help me.” I am disturbed by the pangs of anger, covetousness, gloominess, and driven to disturb the peaceful state in which I was, and which was dear to me: In order that I may not be carried away by raging passion into the bitterness of gall, I must cry out with deep groans: “O God, make speed to save me: O Lord, make haste to help me.” I am tried by being puffed up by accidie, vainglory, and pride, and my mind with subtle thoughts flatters itself somewhat on account of the coldness and carelessness of others: In order that this dangerous suggestion of the enemy may not get the mastery over me, I must pray with all contrition of heart: “O God, make speed to save me: O Lord, make haste to help me.” I have gained the grace of humility and simplicity, and by continually mortifying my spirit have got rid of the swellings of pride: In order that the “foot of pride” may not again “come against me,” and “the hand of the sinner disturb me,”\footnote{Ps. xxxv. (xxxvi.) 12.} and that I may not be more seriously damaged by elation at my success, I must cry with all my might, “O God, make speed to save me: O Lord, make haste to help me.” I am on fire with innumerable and various wanderings of soul and shiftiness of heart, and cannot collect my scattered thoughts, nor can I even pour forth my prayer without interruption and images of vain figures, and the recollection of conversations and actions, and I feel myself tied down by such dryness and barrenness that I feel I cannot give birth to any offspring in the shape of spiritual ideas: In order that it may be vouchsafed to me to be set free from this wretched state of mind, from which I cannot extricate myself by any number of sighs and groans, I must full surely cry out: “O God, make speed to save me: O Lord, make haste to help me.” Again, I feel that by the visitation of the Holy Spirit I have gained purpose of soul, steadfastness of thought, keenness of heart, together with an ineffable joy and transport of mind, and in the exuberance of spiritual feelings I have perceived by a sudden illumination from the Lord an abounding revelation of most holy ideas which were formerly altogether hidden from me: In order that it may be vouchsafed to me to linger for a longer time in them I must often and anxiously exclaim: “O God, make speed to save me: O Lord, make haste to help me.” Encompassed by nightly horrors of devils I am agitated, and am disturbed by the appearances of unclean spirits, my very hope of life and salvation is withdrawn by the horror of fear. Flying to the safe refuge of this verse, I will cry out with all my might: “O God, make speed to save me: O Lord, make haste to help me.” Again, when I have been restored by the Lord’s consolation, and, cheered by His coming, feel myself encompassed as if by countless thousands of angels, so that all of a sudden I can venture to seek the conflict and provoke a battle with those whom a while ago I dreaded worse than death, and whose touch or even approach I felt with a shudder both of mind and body: In order that the vigour of this courage may, by God’s grace, continue in me still longer, I must cry out with all my powers: “O God, make speed to save me: O Lord, make haste to help me.” We must then ceaselessly and continuously pour forth the prayer of this verse, in adversity that we may be delivered, in prosperity that we may be preserved and not puffed up. Let the thought of this verse, I tell you, be conned over in your breast without ceasing. Whatever work you are doing, or office you are holding, or journey you are going, do not cease to chant this. When you are going to bed, or eating, and in the last necessities of nature, think on this. This thought in your heart maybe to you a saving formula, and not only keep you unharmed by all attacks of devils, but also purify you from all faults and earthly stains, and lead you to that invisible and celestial contemplation, and carry you on to that ineffable glow of prayer, of which so few have any experience. Let sleep come upon you still considering this verse, till having been moulded by the constant use of it, you grow accustomed to repeat it even in your sleep. When you wake let it be the first thing to come into your mind, let it anticipate all your waking thoughts, let it when you rise from your bed send you down on your knees, and thence send you forth to all your work and business, and let it follow you about all day long. This you should think about, according to the Lawgiver’s charge, “at home and walking forth on a journey,”\footnote{Deut. vi. 7.} sleeping and waking. This you should write on the threshold and door of your mouth, this you should place on the walls of your house and in the recesses of your heart so that when you fall on your knees in prayer this may be your chant as you kneel, and when you rise up from it to go forth to all the necessary business of life it may be your constant prayer as you stand.

\xsection{Chapter XI. Of the perfection of prayer to which we can rise by the system described.}

Of the perfection of prayer to which we can rise by the system described.

\textsc{This}, this is the formula which the mind should unceasingly cling to until, strengthened by the constant use of it and by continual meditation, it casts off and rejects the rich and full material of all manner of thoughts and restricts itself to the poverty of this one verse, and so arrives with ready ease at that beatitude of the gospel, which holds the first place among the other beatitudes: for He says “Blessed are the poor in spirit, for theirs is the kingdom of heaven.”\footnote{S. Matt. v. 3.} And so one who becomes grandly poor by a poverty of this sort will fulfil this saying of the prophet: “The poor and needy shall praise the name of the Lord.”\footnote{Ps. lxxiii. (lxxiv.) 21.} And indeed what greater or holier poverty can there be than that of one who knowing that he has no defence and no strength of his own, asks for daily help from another’s bounty, and as he is aware that every single moment his life and substance depend on Divine assistance, professes himself not without reason the Lord’s bedesman, and cries to Him daily in prayer: “But I am poor and needy: the Lord helpeth me.”\footnote{Ps. xxxix. (xl.) 17 (LXX.).} And so by the illumination of God Himself he mounts to that manifold knowledge of Him and begins henceforward to be nourished on sublimer and still more sacred mysteries, in accordance with these words of the prophet: “The high hills are a refuge for the stags, the rocks for the hedgehogs,”\footnote{Ps. ciii. (civ.) 18.} which is very fairly applied in the sense we have given, because whosoever continues in simplicity and innocence is not injurious or offensive to any one, but being content with his own simple condition endeavours simply to defend himself from being spoiled by his foes, and becomes a sort of spiritual hedgehog and is protected by the continual shield of that rock of the gospel, i.e., being sheltered by the recollection of the Lord’s passion and by ceaseless meditation on the verse given above he escapes the snares of his opposing enemies. And of these spiritual hedgehogs we read in Proverbs as follows: “And the hedgehogs are a feeble folk, who have made their homes in the rocks.”\footnote{Prov. xxx. 26 (LXX.).} And indeed what is feebler than a Christian, what is weaker than a monk, who is not only not permitted any vengeance for wrongs done to him but is actually not allowed to suffer even a slight and silent feeling of irritation to spring up within? But whoever advances from this condition and not only secures the simplicity of innocence, but is also shielded by the virtue of discretion, becomes an exterminator of deadly serpents, and has Satan crushed beneath his feet, and by his quickness of mind answers to the figure of the reasonable stag, this man will feed on the mountains of the prophets and Apostles, i.e., on their highest and loftiest mysteries. And thriving on this pasture continually, he will take in to himself all the thoughts of the Psalms and will begin to sing them in such a way that he will utter them with the deepest emotion of heart not as if they were the compositions of the Psalmist, but rather as if they were his own utterances and his very own prayer; and will certainly take them as aimed at himself, and will recognize that their words were not only fulfilled formerly by or in the person of the prophet, but that they are fulfilled and carried out daily in his own case. For then the Holy Scriptures lie open to us with greater clearness and as it were their very veins and marrow are exposed, when our experience not only perceives but actually anticipates their meaning, and the sense of the words is revealed to us not by an exposition of them but by practical proof. For if we have experience of the very state of mind in which each Psalm was sung and written, we become like their authors and anticipate the meaning rather than follow it, i.e., gathering the force of the words before we really know them, we remember what has happened to us, and what is happening in daily assaults when the thoughts of them come over us, and while we sing them we call to mind all that our carelessness has brought upon us, or our earnestness has secured, or Divine Providence has granted or the promptings of the foe have deprived us of, or slippery and subtle forgetfulness has carried off, or human weakness has brought about, or thoughtless ignorance has cheated us of. For all these feelings we find expressed in the Psalms so that by seeing whatever happens as in a very clear mirror we understand it better, and so instructed by our feelings as our teachers we lay hold of it as something not merely heard but actually seen, and, as if it were not committed to memory, but implanted in the very nature of things, we are affected from the very bottom of the heart, so that we get at its meaning not by reading the text but by experience anticipating it. And so our mind will reach that incorruptible prayer to which in our former treatise, as the Lord vouchsafed to grant, the scheme of our Conference mounted, and this is not merely not engaged in gazing on any image, but is actually distinguished by the use of no words or utterances; but with the purpose of the mind all on fire, is produced through ecstasy of heart by some unaccountable keenness of spirit, and the mind being thus affected without the aid of the senses or any visible material pours it forth to God with groanings and sighs that cannot be uttered.

\xsection{Chapter XII. A question as to how spiritual thoughts can be retained without losing them.}

A question as to how spiritual thoughts can be retained without losing them.

\textsc{Germanus}: We think that you have described to us not only the system of this spiritual discipline for which we asked, but perfection itself; and this with great clearness and openness. For what can be more perfect and sublime than for the recollection of God to be embraced in so brief a meditation, and for it, dwelling on a single verse, to escape from all the limitations of things visible, and to comprise in one short word the thoughts of all our prayers. And therefore we beg you to explain to us one thing which still remains; viz., how we can keep firm hold of this verse which you have given us as a formula, in such a way that, as we have been by God’s grace set free from the trifles of worldly thoughts, so we may also keep a steady grasp on all spiritual ones.

\xsection{Chapter XIII. On the lightness of thoughts.}

On the lightness of thoughts.

\textsc{For} when the mind has taken in the meaning of a passage in any Psalm, this insensibly slips away from it, and ignorantly and thoughtlessly it passes on to a text of some other Scripture. And when it has begun to consider this with itself, while it is still not thoroughly explored, the recollection of some other passage springs up, and shuts out the consideration of the former subject. From this too it is transferred to some other, by the entrance of some fresh consideration, and the soul always turns about from Psalm to Psalm and jumps from a passage in the Gospels to read one in the Epistles, and from this passes on to the prophetic writings, and thence is carried to some spiritual history, and so it wanders about vaguely and uncertainly through the whole body of the Scriptures, unable, as it may choose, either to reject or keep hold of anything, or to finish anything by fully considering and examining it, and so becomes only a toucher or taster of spiritual meanings, not an author and possessor of them. And so the mind, as it is always light and wandering, is distracted even in time of service by all sorts of things, as if it were intoxicated, and does not perform any office properly. For instance, while it is praying, it is recalling some Psalm or passage of Scripture. While it is chanting, it is thinking about something else besides what the text of the Psalm itself contains. When it repeats a passage of Scripture, it is thinking about something that has to be done, or remembering something that has been done. And in this way it takes in and rejects nothing in a disciplined and proper way, and seems to be driven about by random incursions, without the power either of retaining what it likes or lingering over it. It is then well for us before everything else to know how we can properly perform these spiritual offices, and keep firm hold of this particular verse which you have given us as a formula, so that the rise and fall of our feelings may not be in a state of fluctuation from their own lightness, but may lie under our own control.

\xsection{Chapter XIV. The answer how we can gain stability of heart or of thoughts.}

The answer how we can gain stability of heart or of thoughts.

\textsc{Isaac}: Although, in our former discussion on the character of prayer, enough was, as I think, said on this subject, yet as you want it repeated to you again, I will give you a brief instruction on steadfastness of heart. There are three things which make a shifting heart steadfast, watchings, meditation, and prayer, diligence in which and constant attention will produce steadfast firmness of mind. But this cannot be secured in any other way unless all cares and anxieties of this present life have been first got rid of by indefatigable persistence in work dedicated not to covetousness but to the sacred uses of the monastery, that we may thus be able to fulfil the Apostle’s command: “Pray without ceasing.”\footnote{1 Thess. v. 17.} For he prays \emph{too little}, who is accustomed only to pray at the times when he bends his knees. But he \emph{never} prays, who even while on his bended knees is distracted by all kinds of wanderings of heart. And therefore what we would be found when at our prayers, that we ought to be before the time of prayer. For at the time of its prayers the mind cannot help being affected by its previous condition, and while it is praying, will be either transported to things heavenly, or dragged down to earthly things by those thoughts in which it had been lingering before prayer.

Thus far did Abbot Isaac carry on his Second Conference on the character of Prayer to us astonished hearers; whose instruction on the consideration of that verse quoted above (which he gave as a sort of outline for beginners to hold) we greatly admired, and wished to follow very closely, as we fancied that it would be a short and easy method; but we have found it even harder to observe than that system of ours by which we used formerly to wander here and there in varied meditations through the whole body of the Scriptures without being tied by any chains of perseverance. It is then certain that no one is kept away from perfection of heart by not being able to read, nor is rustic simplicity any hindrance to the possession of purity of heart and mind, which lies close at hand for all, if only they will by constant meditation on this verse keep the thoughts of the mind safe and sound towards God.
